\section{Introduction}

When a language model generates a TriviaQA answer, it leaves two independent traces of its confidence: how consistently its stochastic outputs cluster into the same semantic equivalence class (Semantic Entropy; SE), and how confidently it predicts each individual token on a deterministic greedy decode pass (minimum log-probability; min\_logprob). On TriviaQA dev with Llama-3.1-8B, the Pearson correlation between these two signals is $|r| = 0.026$ --- empirically near-zero, far below the collinearity threshold of 0.70 and the reparameterization boundary of 0.85. This near-orthogonality, while theoretically motivated by the algebraic structure of each signal, had not been directly measured before in the open-weight LLM hallucination detection literature.

The practical stakes are high. Large language models are increasingly deployed in high-stakes domains --- healthcare, legal analysis, knowledge retrieval --- where hallucinated outputs carry real consequences. Uncertainty quantification (UQ) is the primary mechanism for flagging untrustworthy predictions: a reliable UQ signal allows downstream systems to abstain, escalate, or request human review rather than acting on a hallucinated answer. The need for accurate, calibrated UQ signals for open-weight LLMs has motivated a large and growing body of work \citep{guo2017calibration, xiong2023llms, kang2025uncertainty}.

Two dominant signal families have emerged. First, token log-probability signals --- minimum log-probability, mean log-probability, sequence perplexity --- measure the model's token-level confidence from a single greedy decoding pass and require only standard generation API access \citep{malinin2021uncertainty}. On TriviaQA with Llama-3.1-8B, min\_logprob achieves AUROC $\sim$0.825 (our prior internal pipeline experiment h-m1, under identical experimental conditions --- not published). Second, consistency-based signals --- most notably Semantic Entropy (SE) \citep{kuhn2023semantic} --- measure the semantic diversity of $N$ stochastic samples, capturing uncertainty at the semantic, sentence level rather than the token level. SE has shown strong standalone AUROC performance on factual QA benchmarks \citep{kuhn2023semantic, bouchard2025uncertainty}.

The natural next step --- combining SE and min\_logprob in an ensemble --- rests on a critical assumption: that these signals are statistically independent, i.e., that SE captures uncertainty dimensions not already captured by min\_logprob. Without measuring this independence, an ensemble might merely recombine redundant information. Remarkably, this assumption had not been directly tested. The closest related measurement --- first-token confidence vs.\ SE correlation --- yields $r = 0.54$--$0.76$ on Llama-3.1-8B, TriviaQA \citep{gabriel2026first}, suggesting moderate correlation for first-token signals. Whether min\_logprob (minimum over all tokens in greedy decode, not just the first) exhibits the same correlation with SE was an open question.

\textbf{The key insight} driving our work is that SE and min\_logprob are algebraically distinct operations by construction: SE marginalizes over semantic equivalence class probability masses computed from $N=5$ stochastic samples at temperature 0.7, while min\_logprob extracts the minimum token log-probability from a single greedy decoding path. These operations run on different inference calls, aggregate at different granularities (semantic cluster entropy vs.\ token-level minimum), and access different information about the model's output distribution. This algebraic separation should manifest as statistical independence --- and it does: Pearson $|r| = 0.026$ at $N=2500$.

Our contributions are:

\begin{enumerate}
\item \textbf{Empirical independence measurement:} We directly measure Pearson $|r|(\text{SE}_{N5}, \text{min\_logprob}) = 0.026$ on TriviaQA dev with Llama-3.1-8B-Instruct ($N=2500$), confirming near-orthogonality and ruling out SE as a reparameterization of min\_logprob.

\item \textbf{Circularity-controlled evaluation protocol:} Cross-model LM-judge design (Qwen-2.5-7B evaluating Llama-3.1-8B outputs) yields Spearman $\rho(\text{SE}, \text{judge}) = -0.026$ --- a replicable best practice for bias-robust UQ evaluation. We note that this value shares the same rounded magnitude as the Pearson $|r|$ above; both are computed on different variable pairs via different correlation measures, and the coincidence is a feature of this evaluation's near-zero correlation structure.

\item \textbf{Well-powered null result for linear conditional contribution:} At $N=2500$, SE adds no meaningful conditional predictive power beyond min\_logprob (partial $R^2=0.0005$, LRT $p=0.442$). The pre-registered gate evaluation classified this as \texttt{EXPLORE\_N10} (gate\_pass: false), indicating the pipeline recommends extending to $N=10$ stochastic samples per prompt to resolve marginal uncertainty. We explain the power analysis basis for characterizing this as a powered null in Section~\ref{sec:results} and Section~\ref{sec:discussion}, while disclosing the gate outcome honestly.

\item \textbf{End-to-end reproducible pipeline:} A checkpoint-aware, GPU-efficient implementation that completed the full $N=2500$ evaluation on TriviaQA dev --- demonstrating production-quality execution at scale for this hypothesis.
\end{enumerate}

The rest of this paper is organized as follows. Section~\ref{sec:related} positions our work relative to prior SE, log-probability, and LM-judge evaluation literature. Section~\ref{sec:method} describes our methodology. Section~\ref{sec:setup} details the experimental setup. Section~\ref{sec:results} presents results. Section~\ref{sec:discussion} discusses implications and limitations. Section~\ref{sec:conclusion} concludes.
